\documentclass[11pt]{article}

\usepackage[margin=1in]{geometry}
\usepackage[T1]{fontenc}
\usepackage[utf8]{inputenc}
\usepackage{graphicx}
\usepackage{amsmath}
\usepackage{amssymb}
\usepackage{bm}
\usepackage{float}
\usepackage{hyperref}
\usepackage{xcolor}
\usepackage{setspace}
\usepackage{enumitem}
\usepackage{caption}

\singlespacing
\setlength{\parindent}{0pt}
\setlength{\parskip}{0.45em plus 0.1em minus 0.05em}
\graphicspath{{revision_figure_folder/}}

\definecolor{ReviewerSkyBlue}{RGB}{0,150,220}
\definecolor{ManuscriptEditRed}{RGB}{150,0,0}
\definecolor{TodoBrightRed}{RGB}{255,0,0}
\definecolor{StatusGray}{RGB}{90,90,90}

\newcommand{\mschange}[1]{%
  \par\smallskip
  \noindent{\color{ManuscriptEditRed}#1}\par
  \smallskip
}
\newcommand{\statusnote}[1]{\noindent{\color{StatusGray}\textit{#1}}\par\smallskip}
\newenvironment{reviewercommentenv}{\par\smallskip\noindent\color{ReviewerSkyBlue}\itshape\ignorespaces}{\par\normalfont\color{black}\smallskip}
\newcommand{\reviewercomment}[1]{\begin{reviewercommentenv}#1\end{reviewercommentenv}}
\newcommand{\responselabel}{\par\smallskip\noindent\color{black}}
\newcommand{\comment}[1]{{\color{StatusGray}\textit{[comment: #1]}}}

\begin{document}

\begin{center}
{\Large \textbf{Response to Reviewers}}\\[0.5em]
{\large Interspecies Interactions Drive Community-Level Selection in Microbial Coalescence}
\end{center}

Dear Editor and Reviewers,

We thank the reviewers for their careful, rigorous, and constructive feedback. The comments substantially improved both the framing and the technical clarity of the manuscript. Below, we respond point by point and summarize the analyses and text revisions undertaken in response to the major concerns.

\textbf{Summary of major revisions.}

In response to the reviewers, we made four broad changes to the manuscript. First, we clarified the central conceptual framing by defining community-level selection operationally as origin-correlated persistence during coalescence, rather than as evidence for one exclusive biochemical mechanism. Second, we refined the connection between model and experiment by distinguishing the phenomenological gLV interaction-strength parameter from the experimentally measured changes in invasion resistance, net interaction intensity, and environmental feedbacks across nutrient conditions. Third, we added new analyses and controls testing classification-boundary sensitivity, additive null expectations, richness and abundance-imbalance effects, pH-mediated feedback, dominant-species effects, alternative interaction structures, assembly filtering, and taxonomic distinctness in natural sample-derived communities. Finally, we revised the main figures, Extended Data figures, Supplementary Information, and captions to make the assumptions, null models, sample sizes, and scope of the conclusions more explicit.

% ============================================
% Response to Reviewer 3
% ============================================
\section*{Response to Reviewer 3}

We thank the reviewer for their rigorous and technically insightful feedback, particularly the geometric analysis of dimensionality effects on similarity metrics. Below we address each point.

% ============================================
\subsection*{R3-1. Clarify L\_1 versus L\_2 normalization in the similarity metric (P1)}

\reviewercomment{With respect to the similarity metrics used by the authors (the dot products $x_A \cdot x_C$ and $x_B \cdot x_C$), I believe that their interpretation is less straightforward than it would appear at first glance. As a first small comment, the authors state that $x$ are ``normalized abundance vectors,'' and Fig. 1B clearly suggests that these are $L_1$ normalized, that is, their elements sum to 1 (i.e., they represent species' relative abundances). But in the Supplementary Methods, the authors clarify that $x$ are $L_2$ normalized, i.e., it is the sum of squares of their elements that is 1. I think making this more clear in the main text would help, since it fundamentally changes what one should expect for the dot products under different hypotheses.}

\responselabel
We thank the reviewer for identifying this ambiguity and apologize for the confusion. To clarify, $\vec{x}$ denotes a community composition vector over ASVs, and the similarity between the coalesced community C and a parent P was computed as $\mathrm{Sim}(C,P) = (\vec{x}_C \cdot \vec{x}_P)/(\|\vec{x}_C\|_2\|\vec{x}_P\|_2)$. Because cosine similarity is invariant to multiplying either vector by a positive scalar, converting each sample from an ASV-abundance vector to a relative-abundance vector before this step gives the same similarity score; the metric depends on the vector direction, not total abundance. We have corrected the Results, Methods, and Supplementary Note~1 wording accordingly. In the experimental analyses, these composition vectors are ASV-abundance vectors derived from the 16S profiles.


Results \S2.1, similarity-score definition: \mschange{``We represent each community by a community composition vector and compute the cosine similarity between the coalesced community C and each parental community.''} The revised equation explicitly divides each dot product by the product of the two Euclidean norms, and the following sentence defines $\vec{x}_A$, $\vec{x}_B$, and $\vec{x}_C$ as community composition vectors.

Methods \S 16S rRNA Sequencing now states: \mschange{``In experimental analyses, the community composition vector was the ASV-abundance vector derived from these 16S profiles.''}

Methods \S Classification of Coalescence Outcomes, at first definition of the similarity metric: \mschange{``Similarity was computed as cosine similarity between community composition vectors (see Supplementary Note~1 for full normalization details).''}

Supplementary Note~1, metric definition paragraph: \mschange{``To characterize coalescence outcomes, we quantified the compositional similarity between the post-coalescence community and each parental community using cosine similarity.''} The note then gives the same equation as the Results and clarifies that sample-wise total-abundance normalization does not change the cosine-similarity score.

% ============================================
\subsection*{R3-2. Dimensionality artifact in similarity metrics (P0)}

\reviewercomment{%
Along these lines, my main concern is that these similarity metrics can behave in somewhat counter-intuitive ways, and some of the results the authors report may be at least partially skewed by simple statistical effects. To make my point as clear as possible, let me propose a simple example. Let's consider two communities A and B, each with N non-overlapping species. Total species' abundances can be represented by vectors $n_A$ and $n_B$:

\[
n_A = ( n_A^{(1)} , n_A^{(2)} , \ldots , n_A^{(N)} , 0 , 0 , \ldots , 0 )
\]
\[
n_B = ( 0 , 0 , \ldots , 0 , n_B^{(1)} , n_B^{(2)} , \ldots , n_B^{(N)} )
\]

Now consider a simple null model (similar to the ``shuffled abundance'' null model proposed by the authors) in which community C, which results from A and B undergoing coalescence, has N species. The abundance vector of C, $n_C$, has N elements that are zero and N that are non-zero. Which elements are set to zero is chosen randomly (though the argument holds in other scenarios, including the ``abundance-weighted random selection'' model formulated by the authors). The non-zero elements of $n_C$ are random samples of the non-zero elements of $n_A$ and $n_B$. This null model aims at representing a ``maximal restructuring'' situation, where species' abundances in community C are completely uncorrelated from their abundances in the parental communities. If we generate $n_A$ and $n_B$ by randomly sampling their non-zero elements uniformly in (0, 1) and we compute the similarity metrics

\[
\mathrm{Sim}(A,C) = n_A \cdot n_C / |n_A| |n_C|
\]
\[
\mathrm{Sim}(B,C) = n_B \cdot n_C / |n_B| |n_C|
\]

For different values of N and different pairs of randomly-generated communities, the ``similarity map'' looks like:

[See attached Reviewer Figure 1]

So despite the model being formulated as a null baseline representing a ``full restructuring'' situation, many cases would be flagged as ``dominance'' or ``mixing.'' The same is true if species' abundances in A and B are sampled differently, for instance if instead of sampling them uniformly ($n \sim U(0,1)$) we used a more uneven distribution (such as $n \sim 10^{U(-3,0)}$) we would find:

[See attached Reviewer Figure 2]

Notably, a somewhat similar effect also occurs under a much simpler null model: if the community C was simply the superposition of communities A and B (i.e., if the two parental communities mixed perfectly without interacting), we would have $n_C = n_A + n_B$. In this case:

[See attached Reviewer Figure 3]

The key takeaway here is that, when using these similarity metrics and the author's ``Dominance-Mixing-Restructuring'' classification, one could expect to observe a decline in the ``Dominance'' fraction as community diversity increases due to a simple statistical effect: positive unit vectors are more likely to be close to orthogonal in lower dimensions than they are in higher dimensions.%
}

\begin{figure}[H]
\centering
\includegraphics[width=0.88\textwidth,height=0.78\textheight,keepaspectratio]{Reviewer3_attachment_figures.pdf}
\caption*{\textbf{Reviewer attachment.}}
\end{figure}

\responselabel
We thank the reviewer for raising this important point. We reproduced the reviewer's toy null models plus the corresponding skewed-abundance additive case (random restructuring with $U(0,1)$, random restructuring with $10^{U(-3,0)}$, simple additive mixing $n_C = n_A + n_B$ with abundances drawn from $U(0,1)$, and simple additive mixing $n_C = n_A + n_B$ with abundances drawn from $10^{U(-3,0)}$) using the exact same L$_2$/cosine analysis used in our manuscript (Response Fig.~R3-2A). We note that one mismatch in the reviewer-attached visualization is that the Restructuring boundary appears to use a retention-magnitude cutoff of $r < 1/2$, whereas our pipeline uses $r \leq 1/\sqrt{2}$ (equivalently, $r^2 \leq 1/2$).

Using the same classification pipeline as in the manuscript, the reproduced results support the reviewer's main concern in a specific regime. The random-restructuring null models become increasingly classified as Restructuring as $N$ grows (uniform case: $43.4\%, 64.9\%, 78.5\%, 88.7\%$ for $N = 2, 4, 6, 8$), and the uniform additive null model concentrates on the symmetric Mixture axis and is classified as Mixture in $80.5\%, 97.4\%, 99.1\%, 99.7\%$ of cases. However, the skewed-abundance random-restructuring null remains high-retention but can fall off the symmetry line and be classified as high Dominance (Dominance: $73.9\%, 50.0\%, 39.7\%, 30.6\%$ for $N = 2, 4, 6, 8$). This confirms that skewed abundance distributions can generate apparent Dominance under simple null models, motivating the event-matched additive-null analysis below.

We have revised the Methods and Supplementary Note~1 to state this boundary explicitly, and Response Fig.~R3-2A shows both the manuscript boundary and the alternative visual boundary at $r = 1/2$.

Methods, ``Classification of Coalescence Outcomes,'' now states that \mschange{``The Restructuring boundary is defined by $r^2 \leq 0.5$, equivalently a retention radius $r \leq 1/\sqrt{2}$ in the similarity map.''} Supplementary Note~1, ``Outcome classification,'' now states that \mschange{``Restructuring occurs when $r^2 \leq 0.5$, equivalently when the retention radius satisfies $r \leq 1/\sqrt{2}$ in the two-parental-community similarity map,''} and \mschange{``the retention threshold $r^2 = 0.5$ (radius $r = 1/\sqrt{2}$, not $r = 1/2$) marks the point at which the best linear combination of the two-parental-community composition vectors accounts for half of the coalesced community's composition in the normalized similarity space.''}

\begin{figure}[H]
\centering
\includegraphics[width=0.95\textwidth]{Fig_R3_2_reviewer_reproduction_L2.pdf}
\caption*{\textbf{Response Fig.~R3-2A.} \textbf{Reviewer toy nulls analyzed with the manuscript classifier.} Rows: \textbf{A}, random restructuring with $U(0,1)$ abundances; \textbf{B}, random restructuring with $10^{U(-3,0)}$ abundances; \textbf{C}, additive mixing with $U(0,1)$ abundances; \textbf{D}, additive mixing with $10^{U(-3,0)}$ abundances. Scatter panels show the reviewer's visual boundary ($r=1/2$, gray) and the manuscript boundaries (retention radius $r=1/\sqrt{2}$ and asymmetry class boundaries, black); bars show manuscript-classifier fractions.}
\end{figure}

\reviewercomment{%
This effect is partially accounted for in the manuscript; specifically, Extended Data Fig.~3 shows that the empirical distribution of Dominance values has significantly higher mean than its counterparts generated by the null models. However, while the authors' claim is correct (``skewed abundance distributions alone do not fully account for the observed asymmetric outcomes''), this does not mean that the empirically quantified Dominance values aren't biased by the statistical effects described above. Many cases flagged as ``Dominance'' may in actuality be compatible with a simple null model, particularly for the communities assembled from the less diverse species pools.%
}

\responselabel
The prior analysis demonstrates that a simple additive null model with skewed abundance can be classified as Dominance because of a simple statistical effect. Specifically, this occurs when the two parental-community abundance vectors have substantially different norms. For a simple additive null model, $n_A+n_B$ is closer to the larger-norm parental community. For non-overlapping parental-community vectors, this effect is direct: $\mathrm{Sim}(A,A+B)$ and $\mathrm{Sim}(B,A+B)$ are proportional to $\|n_A\|_2$ and $\|n_B\|_2$, respectively. We agree that this could particularly matter when abundance is strongly uneven, because vector norms become sensitive to concentration of abundance among a small number of taxa.

However, the norm-imbalance effect is much weaker in our experimental parent pairs than in the skewed toy examples. In our experimental coalescence pairs, the fold differences in Euclidean norms between the two parental raw 16S read-count vectors were $1.85 \pm 0.95$ (mean $\pm$ SD; median $1.58$), with per-medium means ranging from $1.55 \pm 0.51$ to $2.05 \pm 1.18$. These values are far smaller than those in the skewed-abundance simple additive null model, where the fold differences were $18.2 \pm 39.4$ and $7.1 \pm 14.7$ for $N=2$ and $N=4$, respectively (Response Fig.~R3-2B).

To directly test whether this effect could account for the experimental classifications, we constructed a simple additive null model from the actual raw 16S read-count vectors of the same two parental communities in each experimental coalescence pair (Response Fig.~R3-2B). The simple additive null model was classified as Mixture in most events in each nutrient condition: 84/90 in Nutr$-$ (93.3\%), 66/83 in Base (79.5\%), and 69/90 in Nutr$+$ (76.7\%). Thus, the effect identified by the reviewer is present in principle and detectable in some experimental parent pairs, but it is not sufficient to explain the observed Dominance tendency.

\begin{figure}[H]
\centering
\includegraphics[width=0.92\textwidth]{Fig_R3_2_parent_norm_asymmetry.pdf}
\caption*{\textbf{Response Fig.~R3-2B.} \textbf{Parental norm imbalance and simple additive null model classifications.} \textbf{A}, Raw 16S read-count L$_2$ norms for parental-community abundance vectors. \textbf{B}, Paired parental L$_2$ norm fold differences, with skewed-abundance simple additive null models shown for scale; dashed line, $1+\sqrt{2}$ fold difference corresponding to PDI $<0.25$ or $>0.75$ for the simple additive null model. \textbf{C}, Simple additive null model class fractions; labels show Mixture counts.}
\end{figure}

\reviewercomment{%
In my view, it would be necessary to compare each coalescence community with a corresponding null model constructed specifically for each of the pairs of parental communities, such as the simple additive model $n_{C,\mathrm{null}} = n_A + n_B$. Comparing the experimental coalesced community with such a null model should be done on a case-by-case basis, not just by comparing distributions of Dominance indexes as in Extended Data Fig.~3.%
}

\responselabel
We thank the reviewer for this suggestion. We compared each observed Base-medium coalescence community with the corresponding simple additive null model, constructed from the same two parental communities using raw 16S read-count vectors. As described above, this null model was usually classified as Mixture, although it produced Dominance in a subset of events. Critically, the observed Base-medium outcomes were shifted further toward direction-independent parental asymmetry, $y = |2\mathrm{PDI}-1|$, than their matched additive-null outcomes (77/83 events; one-sided paired Wilcoxon test, $p = 5.9 \times 10^{-14}$). Thus, the event-matched comparison shows that the simple additive model can contribute to apparent asymmetry in some cases but does not by itself explain the observed Dominance tendency. We replaced Extended Data Fig.~3 with this event-matched additive-null analysis and described the associated norm-imbalance effect in Supplementary Note~1.

\begin{figure}[H]
\centering
\includegraphics[width=0.98\textwidth]{Fig_R3_2_additive_null_comparison.pdf}
\caption*{\textbf{Response Fig.~R3-2C.} \textbf{Base-medium raw-count observations versus the event-matched simple additive null model.} \textbf{A}, Raw-count similarity map: orange open points, simple additive null model outcomes ($n_{C,\mathrm{null}}=n_A+n_B$); blue filled points, observed coalesced communities; red arrows, five example matched null-to-observation pairs. \textbf{B}, Paired shifts in direction-independent parental asymmetry, $y = |2\mathrm{PDI}-1|$, from additive null to observation for all 83 Base events; gray lines connect matched pairs and the dashed line marks the $y=0.5$ asymmetry boundary. \textbf{C}, Class transitions from additive-null outcomes (columns) to observations (rows).}
\end{figure}

% ============================================
\subsection*{R3-3. Interaction strength, diversity, and Dominance frequency (P0)}

\reviewercomment{In section 2.3 of the Results, the authors examine how varying the mean interaction (competition) strength between species in a gLV model affects coalescence outcomes. The main conclusion here is that increasing competition strength shifts the distribution of coalescence outcomes from all ``Mixing'' to mostly ``Dominance'' and ``Restructuring.'' However, it is not clear to what extent this phenomenon can be explained by simple statistical effects like the ones described above. Increasing the mean interaction strength (mu) can be expected to decrease the number of species in the communities, and, as mentioned above, reducing N generically causes an increase in the frequency of ``Dominance'' outcomes. It seems important to test whether the increase of ``Dominance'' outcomes in the model (and in the experiments) as mu grows is substantially different from the expected increase in the null models as N decreases.}

\responselabel
As shown in the R3-2 analyses, in our experiment the richness and norm-imbalance effect identified by the reviewer can occur, but it is not sufficient to explain the observed Dominance increase. For the experimental data, the simple additive null model shows that paired parental raw 16S read-count vector norm differences were relatively small and that the simple additive null model was usually classified as Mixture rather than Dominance (Response Fig.~R3-2B). Thus, this statistical effect contributes to the null expectation in some events but does not account for the experimental Dominance pattern.

In the gLV simulations, increasing the interaction-strength parameter reduces final richness and creates greater unevenness in parental-community norms, so Dominance outcomes in the simple additive null model become more frequent because of this richness- and norm-imbalance effect (Response Fig.~R3-3A). However, this effect is insufficient to explain Dominance in the full simulated coalescence outcomes. At the three representative interaction-strength parameter values used in the main figure, the simple additive null model gives $0.0\%$, $5.7\%$, and $16.2\%$ Dominance at $\mu = 0.3$, $0.6$, and $0.8$, respectively, whereas the full gLV coalescence simulations give $20.4\%$, $61.0\%$, and $69.5\%$ Dominance. Thus, richness- and norm-imbalance effects contribute to the null expectation at high $\mu$, but they do not explain the simulated Dominance transition.

We have added this analysis to Supplementary Note~5, in the paragraph on abundance-skew and richness-dependent geometric effects, and show the full analysis in Supplementary Fig.~38.

\begin{figure}[H]
\centering
\includegraphics[width=0.90\textwidth]{Fig_R3_3_sim_parent_norm_asymmetry.pdf}
\caption*{\textbf{Response Fig.~R3-3A.} \textbf{Simulation parental-community composition-vector norm imbalance and simple additive null model classifications across interaction strength.} \textbf{A}, Euclidean norms of assembled parental-community composition vectors across $\mu$; line shows the median and shading shows the 25th--75th percentile range. \textbf{B}, Fold difference in Euclidean norm between the two parental communities for each simulated coalescence event; line shows the median and shading shows the 25th--75th percentile range, and dotted line shows the 95th percentile. The dashed horizontal line marks the fold difference $1+\sqrt{2}$ corresponding to PDI $<0.25$ or $>0.75$ for the simple additive null model. \textbf{C}, Classification fractions for the simple additive null model computed from the same paired parental-community composition vectors at each $\mu$.}
\end{figure}

\reviewercomment{It would also be useful to see how richness (in the final communities, not in the inocula) changes with mu in the model, as well as across the different media (Base, Nutr-, Nutr+) used in the experiment. From Supplementary Figs. 10-12, it seems like communities derived from natural samples do exhibit substantial diversity differences in the Base and Nutr+ media with respect to the Nutr- medium, but I could not find an equivalent representation for the first experiment.}

\responselabel
We agree that showing the final-community richness directly is useful. In the simulations, final richness decreases with $\mu$ for both post-assembly sub-communities and coalesced communities (Response Fig.~R3-3B, panel A). In the experiment, final coalesced-community richness also differs across media, with mean richness of 13.0, 9.9, and 8.5 ASVs in Nutr$-$, Base, and Nutr$+$, respectively (Response Fig.~R3-3B, panel B). Thus, richness differences are real and are now shown directly. However, as demonstrated by the prior analyses, the associated statistical/geometric effect does not explain the Dominance patterns: it does not account for the experimental event-matched additive-null result (Response Fig.~R3-2B,C) or for the simulation trend (Response Fig.~R3-3A).

\begin{figure}[H]
\centering
\includegraphics[width=0.92\textwidth]{Fig_R3_3_richness_summary.pdf}
\caption*{\textbf{Response Fig.~R3-3B.} \textbf{Final richness across interaction strength and nutrient conditions.} \textbf{A}, Mean final richness of simulated post-assembly sub-communities and coalesced communities across $\mu$; error bars show s.e.m. \textbf{B}, Final coalesced-community richness (ASVs above 0.1\% threshold) across nutrient conditions and initial pool sizes in the synthetic-community experiment.}
\end{figure}

% ============================================
\subsection*{R3-4. gLV model excludes facilitation}

\reviewercomment{With respect to the model, the gLV framework is restricted to competitive interactions ($\alpha_{ij} > 0$). While this is reasonable from a modeling perspective because it prevents situations of ``unchecked growth,'' it also limits the range of ecological scenarios that the model can represent. Previous works have shown that facilitative interactions (particularly in the form of metabolic cross-feeding) are an important determinant of correlated species' fates during coalescence.}

\responselabel
We thank the reviewer for raising this important and valid limitation. We treat competition as the major source of interaction in this experimental system and in our modeling, while recognizing that facilitation can be important in other microbial coalescence settings. This approximation is supported by the 12-isolate pairwise coculture assay, where focal ASV yields were usually lower in coculture than in monoculture: 79/84/93\,\% of species-in-pair observations fell below the monoculture expectation in Nutr$-$/Base/Nutr$+$ (Response Fig.~R3-4a).

We agree that the competition-only gLV model cannot capture facilitative interactions such as metabolic cross-feeding. To consider the broader ecological context, we ran the same coalescence analysis for two model extensions: (1) facilitative tails that allow positive ecological interactions (negative $\alpha_{ij}$ under our sign convention), and (2) reciprocal sign/correlation structure between coefficients $\alpha_{ij}$ and $\alpha_{ji}$. The simulation definitions and full parameter sweeps are now described in Supplementary Note~3 and Supplementary Figs.~39--40.

In the facilitative-tail extension, we preserved the mean interaction coefficient and added density-dependent self-limitation to prevent runaway growth (Response Fig.~R3-4b). Dominance still rises with $\mu$ at every facilitative-tail strength tested. At $\mu=0.30$, increasing the facilitative tail shifts outcomes away from Mixture toward both Dominance and Restructuring (Dom/Mix/Res $=18/73/9\,\%$ at $f=0$ and $43/33/24\,\%$ at $f=0.8$). At $\mu=0.80$, Dominance remains high across the facilitative tail (70--76\%). Thus, facilitative interactions can change the balance between Mixture and Restructuring in some regimes, but the core claim, the monotonic rise of Dominance with stronger interaction ($\mu$), remains largely valid in this model extension.

We also varied reciprocal pair-coupling structure, using antisymmetric exploitation for $p<0$ and symmetric competition for $p>0$ (Response Figs.~R3-4c,d). This analysis shows that reciprocal coefficient correlation does not affect our core claim: Dominance still increases with stronger interaction ($\mu$) when reciprocal coefficients are anticorrelated or positively correlated.

We clarified this scope in the Supplementary Methods: \mschange{``Because we restrict $\alpha_{ij} \geq 0$, the model captures only competitive (growth-inhibiting) interactions; facilitative interactions such as cross-feeding are not represented.''} The Discussion now states the corresponding limitation: \mschange{``Additionally, both our experimental system and theoretical model focus on a competition-dominated regime; systems with substantial mutualism or facilitation may exhibit qualitatively different dynamics.''}

\begin{figure}[H]
\centering
\includegraphics[width=\textwidth]{R3_4_experiment.pdf}
\caption*{\textbf{Response Fig.~R3-4a.} \textbf{Pairwise coculture data show widespread growth-inhibiting effects, especially in Base and Nutr$+$.}
\textbf{(A--C)} Per-medium scatter of each focal ASV's coculture CFU ($C_i$) versus its monoculture CFU ($M_i$) from the 12-isolate pairwise invasion assay. Each point is one ASV in one unordered pair, averaged across the two invasion directions when both were available; the dashed line marks $C_i=M_i$. Points below the diagonal indicate direct suppression of that ASV in coculture relative to monoculture.
\textbf{(D)} Bar summary of the fraction of species-in-pair observations below the monoculture expectation: 79/84/93\,\% for Nutr$-$/Base/Nutr$+$. These fractions are descriptive summaries; no p-values are shown for this panel.}
\label{fig:R3-4a}
\end{figure}

\begin{figure}[H]
\centering
\includegraphics[width=\textwidth]{R3_4_mixed_sign_higher_order.pdf}
\caption*{\textbf{Response Fig.~R3-4b.} \textbf{Allowing a facilitative tail can increase Restructuring but does not remove the rise of Dominance with $\mu$.}
Off-diagonal coefficients are drawn iid from $\alpha_{ij}\sim U[-f\mu,(2+f)\mu]$, preserving $E[\alpha_{ij}]=\mu$ while increasing the expected facilitative fraction $P(\alpha_{ij}<0)$ from 0 to 22.2\%. Top insets show the analytic sampling density of $\alpha_{ij}/\mu$ (red: facilitative support; gray: nonnegative support), drawn from the formula rather than from simulated matrices. Dynamics include a stabilizing density-dependent self-limitation term, $\dot n_i=n_i(1-(A n)_i-0.1n_i^2)$. Each panel fixes one facilitative-tail strength $f$ and shows Dominance / Mixture / Restructuring fractions at $\mu \in \{0.30,0.60,0.80\}$ (200 pools $\times$ 6 coalescence events per cell). Trends across the simulated parameter grid are descriptive; no p-values are shown for this robustness simulation.}
\label{fig:R3-4b}
\end{figure}

\begin{figure}[H]
\centering
\includegraphics[width=\textwidth]{R3_4_simulation.pdf}
\caption*{\textbf{Response Fig.~R3-4c.} \textbf{Dominance still increases with $\mu$ across reciprocal pair-coupling structures.}
Each panel fixes one $p$ and shows Dominance / Mixture / Restructuring fractions at $\mu \in \{0.30, 0.60, 0.80\}$, ordered from antisymmetric exploitation ($p<0$) through the i.i.d.\ competitive baseline ($p=0$) to symmetric competition ($p>0$). Top insets show the analytic directed-coefficient distribution implied by each $p$ value for $\alpha_{ij}/\mu$ (red: negative coefficients; gray: nonnegative coefficients), drawn from the sampler definition rather than from Monte Carlo samples. Data: 200 pools $\times$ 6 coalescence events per cell. Trends across $\mu$ and $p$ are descriptive; no p-values are shown for this robustness simulation.}
\label{fig:R3-4c}
\end{figure}

\begin{figure}[H]
\centering
\includegraphics[width=\textwidth]{R3_4_pair_coupling_fine.pdf}
\caption*{\textbf{Response Fig.~R3-4d.} \textbf{The finer pair-coupling sweep preserves the $\mu$-dependent Dominance trend.}
Same pair-coupling definition as Response Fig.~R3-4c, but with $p$ swept from $-1$ to $+1$ in increments of 0.2. Top strips show representative analytic directed-coefficient distributions for $p=-1,0,+1$; the sweep itself uses the same sampler for all $p$ values. Each panel fixes $\mu$ and shows stacked Dominance / Mixture / Restructuring fractions across $p$; the black line marks the Dominance fraction. Trends across the simulated $p$ sweep are descriptive; no p-values are shown for this robustness simulation.}
\label{fig:R3-4d}
\end{figure}

\reviewercomment{Would two species which engage in a mutualism be considered ``weak interactors'' or the opposite?}

\responselabel
Within our framework, interaction strength refers to coefficient magnitude and distributional spread (the May-style random-Lotka--Volterra axis; May 1972; Grilli 2017), not interaction sign. Thus, two species in a strong mutualism would be strong interactors, with coefficients of opposite sign from strong competitors. We note that our empirical invasion-resistance assay, however, measures growth-inhibiting effects and would not score mutualistic facilitation as strong competitive exclusion.

To test whether this sign distinction affected the model result, we added weak bidirectional mutualistic pairs to the gLV simulations. Specifically, 0--40\% of unordered species pairs were made bidirectionally mutualistic ($\alpha_{ij},\alpha_{ji}\sim U[-0.2\mu,0]$), while the remaining interactions were left competitive. The $\mu$-dependent Dominance trend persisted: Dominance remained 58.5--70.2\% at $\mu=0.60$ and 69.4--77.5\% at $\mu=0.80$ (Response Fig.~R3-4e; Supplementary Fig.~41). Thus, adding weak mutualistic pairs does not remove the Dominance trend.

This magnitude/spread framing is further supported by the mean-vs-variance gLV sweep $\alpha_{ij}\sim U[m-h,m+h]$ with $m,h \in \{0,0.2,0.4,0.6,0.8\}$: zero coefficient half-width ($h=0$) gave only Mixture outcomes even at the highest mean ($m=0.8$), and across the 25 cells Dominance correlated more strongly with the half-width $h$ than with the mean $m$ (descriptive Pearson $r=0.78$ vs.\ $0.53$; Response Fig.~R3-4f; Supplementary Fig.~42). Mechanistically, heterogeneous coefficients are what let assembly filter species into mutually compatible subgroups whose members face stronger competition from outsiders, the ``cohesion without cooperation'' route to Dominance noted in the Discussion (Tikhonov 2016; May 1972; Grilli 2017). We clarified this mean-vs-variance decomposition in Supplementary Methods (``Lotka--Volterra Simulations'') and Supplementary Note~3, also in response to Reviewer 2's related request for biological interpretation of $\mu$ and the coefficient distributions (R2-6).

\begin{figure}[H]
\centering
\includegraphics[width=\textwidth]{R3_4_mutualistic_pair_fraction.pdf}
\caption*{\textbf{Response Fig.~R3-4e.} \textbf{Weak bidirectional mutualistic pairs do not remove the $\mu$-dependent Dominance trend.}
For a controlled fraction $q \in \{0,0.10,0.20,0.30,0.40\}$ of unordered species pairs, both reciprocal coefficients are sampled from $\alpha_{ij},\alpha_{ji}\sim U[-0.2\mu,0]$, corresponding to weak positive ecological interactions under our sign convention. All remaining off-diagonal coefficients are sampled from the baseline competitive distribution $U[0,2\mu]$. Dynamics include the same stabilizing density-dependent self-limitation term used in Response Fig.~R3-4b, $\dot n_i=n_i(1-(A n)_i-0.1n_i^2)$. Each panel fixes one mutualistic-pair fraction $q$ and shows Dominance / Mixture / Restructuring fractions at $\mu \in \{0.30,0.60,0.80\}$ (200 pools $\times$ 6 coalescence events per cell). Trends across the simulated parameter grid are descriptive; no p-values are shown for this robustness simulation.}
\label{fig:R3-4e}
\end{figure}

\begin{figure}[H]
\centering
\includegraphics[width=0.92\textwidth]{R3_4_mean_variance_grid.pdf}
\caption*{\textbf{Response Fig.~R3-4f.} \textbf{Coefficient variance, not coefficient mean alone, drives Dominance in the mixed-sign gLV sweep.}
Off-diagonal coefficients are sampled from $\alpha_{ij}\sim U[m-h,m+h]$, with mean coefficient $m$ and standard deviation $h/\sqrt{3}$ varied independently. Negative $\alpha_{ij}$ values, corresponding to positive ecological interactions under our sign convention, occur above the dashed $h=m$ line. Dynamics include the same stabilizing density-dependent self-limitation term used in Response Fig.~R3-4b, $\dot n_i=n_i(1-(A n)_i-0.1n_i^2)$. Heatmaps show Dominance, Mixture, Restructuring, and same-parent concordance $|\phi|$ across the 5 $\times$ 5 grid (200 pools $\times$ 6 coalescence events per cell). Trends are descriptive; no p-values are shown for this robustness simulation.}
\label{fig:R3-4f}
\end{figure}

\reviewercomment{This seems particularly relevant given that ``Restructuring'' appears to be much more common in natural sample-derived communities than in the initial bottom-up-assembled communities. Given these discrepancies, and the fact that natural sample-derived communities are still model systems (even if they may be closer to natural microbiomes than communities assembled from defined species pools), I would suggest the claims in section 2.6 of the Results as well as in the Discussion are toned down (e.g., lines 292-294).}

\responselabel
We have changed and toned down our conclusion for natural sample-derived communities in Results \S2.6: \mschange{``These communities showed higher Restructuring fractions, possibly reflecting greater taxonomic diversity and more complex interaction networks. These results suggest that the nutrient-dependent increase in Dominance observed in synthetic consortia is also present in taxonomically richer natural sample-derived communities.''} We also removed the detailed natural-community generality caveat from the Discussion and consolidated it in Supplementary Note~6, which emphasizes that these are laboratory-stabilized communities rather than unfiltered natural ecosystems: \mschange{``Thus, pre-selection could contribute to the convergent coalescence patterns observed between synthetic and natural sample-derived communities, and the natural-community results should be interpreted as evidence that the nutrient-dependent increase in Dominance occurs in taxonomically richer laboratory-stabilized communities, not as evidence for unrestricted generality across unfiltered natural ecosystems.''}

% ============================================
\subsection*{R3-5. ``Interaction strength'' versus ``competition strength''}

\reviewercomment{One last comment along these lines, which is admittedly mostly semantic: the term ``interaction strength'' may be somewhat confusing since, in the context of the gLV model, it refers strictly to competitive interactions. It may be worth considering using different terminology (``competition strength''?)}

\responselabel
We thank the reviewer for raising this important concern about the semantics around the term ``interaction strength''; after revisiting the terminology, we retain this term because it is consistent with the random-interaction and microbial gLV literature that motivates our framework, including May's classic random community model (May 1972) and recent microbial gLV work by Hu et al. (2022, 2025), where interaction strength is used prevalently. The R3-4 analyses support this scope by showing that the empirical coculture data contain frequent growth-inhibiting effects, that the qualitative Dominance trend persists in facilitative-tail and pair-coupling model extensions, and that coefficient variance is required when coefficient mean and variance are separated. We therefore kept the established term while clarifying that, in the baseline gLV model, $\mu$ sets the range of sampled competitive interaction coefficients and controls both their mean and variance. Methods now states: \mschange{``We fixed growth rates to 1 and self-interaction coefficients $\alpha_{ii} = 1$, and drew off-diagonal competition coefficients from a uniform distribution $\mathbb{U}(0, 2\mu)$, so $\mu$ sets the range of sampled competitive interaction coefficients and thereby controls both their mean and variance within this phenomenological model.''} We also clarified in the revised Discussion, quoted above in R3-4, that the model focuses on a competitive-interaction regime.

\clearpage
% ============================================
% Response to Reviewer 2
% ============================================
\section*{Response to Reviewer 2}

We thank the reviewer for their thoughtful and constructive comments, which have significantly improved the conceptual clarity of the manuscript. Below we address each point.

% ============================================
\subsection*{R2-1. Nutrient enrichment versus interaction strength}

\reviewercomment{The central claim of the paper is that increasing nutrient supply strengthens interactions, which in turn drives a shift toward community-level selection. While this is an appealing and intuitive interpretation, it may be overly simplified.

Nutrient enrichment simultaneously affects multiple ecological processes, including:
\begin{itemize}
\item species' carrying capacities
\item metabolic rates
\item environmental modification (e.g.\ pH shifts)
\item the balance of competition and facilitation
\end{itemize}

Indeed, the manuscript itself shows that pH modification by dominant taxa may contribute to winner identity in the nutrient-rich regime, while not serving as a sole explanation for Dominance frequency. This suggests that the observed transition may reflect changes in environmental feedbacks and dominance structure, rather than a simple increase in pairwise interaction strength.

This raises an important conceptual point:

\emph{Is the observed shift in coalescence outcomes driven by increased interaction strength per se, or by other effects of resource enrichment?}

We therefore suggest reframing this result more broadly in terms of changes in \emph{interaction intensity and environmental feedbacks}, rather than interpreting nutrient enrichment as a direct proxy for stronger interactions. This would align well with recent work highlighting the role of environmentally mediated interactions in microbial communities [Goldford et al., 2018, Estrela et al., 2021].}

\responselabel
We agree that this is an important conceptual question to clarify. In our work, the gLV interaction-strength parameter is framed as a phenomenological parameter for ecological coupling within the model, while the nutrient manipulation in Fig.~4 changes nutrient availability and is interpreted through its effects on invasion resistance, net interaction effects, and environmental feedbacks. We therefore use failed pairwise invasion frequency as an operational measure of invasion resistance, in order to approximate how net interaction effects change across nutrient conditions, rather than as a direct measurement of one biochemical mechanism or a literal estimate of every pairwise Lotka--Volterra coefficient. Multiple mechanisms, including resource competition, increased metabolic activity, pH-mediated environmental feedbacks, and changes in the competition--facilitation balance, can contribute to the observed nutrient-dependent change in invasion resistance. This framing is consistent with prior work on nutrient-mediated microbial interactions and environmentally mediated community assembly (Ratzke and Gore, 2018; Ratzke et al., 2020; Goldford et al., 2018; Estrela et al., 2021; Duan et al., 2025).

\reviewercomment{From a consumer--resource perspective, it is also not mathematically self-evident that increasing resource supply monotonically increases effective pairwise interaction strengths. In our recent theoretical work [Duan et al., 2025], we show that when mapping mechanistic consumer--resource dynamics to effective Lotka--Volterra coefficients, the external supply rate can cancel out under broad conditions, provided that resources are not strictly limiting and metabolic strategies remain fixed.

Under such conditions, resource enrichment increases total biomass but does not necessarily increase mean per-capita competition strength. Instead, enrichment restructures the interaction network through changes in equilibrium abundances and environmental feedbacks. This distinction is important, as it suggests that the observed transition may reflect a reorganisation of interaction structure rather than a simple scaling of interaction strength.

A brief discussion of these issues would significantly strengthen the conceptual clarity of the manuscript.}

\responselabel
The apparent discrepancy comes from how the interaction coefficient is defined. The consumer--resource result discussed by Duan, Pawar and colleagues concerns coefficients derived mechanistically from resource dynamics as effective \emph{per-capita} terms in a population-growth equation. In contrast, our gLV coefficients are used as coarse-grained effective parameters in a normalized-abundance model, and our nutrient-gradient experiment uses pairwise invasion assays to quantify empirical invasion resistance. We therefore do not interpret nutrient enrichment as a direct or monotonic increase in those mechanistic \emph{per-capita} Lotka--Volterra coefficients. As the reviewer correctly notes, resource enrichment can increase total biomass without simply scaling the underlying \emph{per-capita} coefficients. Thus, the nutrient-gradient result should be interpreted as an increase in empirical invasion resistance, not as a claim that nutrient supply directly or monotonically scales every effective \emph{per-capita} interaction coefficient.

Results \S2.4 (``Nutrient-dependent interaction intensity and environmental feedbacks recapitulate model predictions'') now reads: \mschange{``Following prior work showing that nutrient concentration can alter microbial competition and that pH-mediated environmental feedbacks can drive bacterial interactions (Hu et~al.\ 2022; Ratzke and Gore 2018; Ratzke et~al.\ 2020; Hu et~al.\ 2025), we conducted additional coalescence experiments by removing or augmenting glucose and urea in the Base medium used in Fig.~1 (Methods). To approximate how net interaction effects change across nutrient conditions, we measured invasion resistance using pairwise invasion assays among the 12 most abundant isolates.''}

% ============================================
\subsection*{R2-2. Alternative explanations for community-level selection}

\reviewercomment{The manuscript interprets dominance of one parental community and correlated persistence of its taxa as evidence for community-level selection. We find this interpretation interesting and potentially insightful, but it is not uniquely diagnostic on its own.

Similar patterns could arise from:
\begin{itemize}
\item shared environmental filtering
\item correlated traits within parental communities
\item environmental modification (e.g.\ pH tolerance)
\end{itemize}

In particular, the nutrient-rich regime appears consistent with strong environmental filtering imposed by a small number of dominant taxa, which may not require invoking selection acting at the level of the community as a unit.

We therefore recommend clarifying the interpretation of ``community-level selection,'' and explicitly discussing alternative mechanisms that could generate similar patterns. This distinction has been discussed in previous work on community coalescence and microbial assembly [Mansour et al., 2018, Rillig et al., 2015].}

\responselabel
We thank the reviewer for raising this important point. We agree that Dominance and correlated persistence should not be interpreted as uniquely identifying a single underlying mechanism. We have therefore revised the manuscript to present Dominance as phenomenological support for community-level selection, rather than as direct evidence for one causal process. We clarify that multiple mechanisms, including shared environmental tolerance, correlated traits, and filtering by pH-modifying dominant taxa, can generate this outcome-level pattern; Dominance alone does not prove one exclusive mechanism.

We made this distinction explicit in Results \S2.1 (``Community-level selection is prevalent in microbial coalescence''): \mschange{``Crucially, Dominance provides phenomenological support for community-level selection: one parental community displaces the other while retaining its internal compositional structure, indicating that its species persist collectively.''}

This point is especially relevant in Nutr$+$, where a few dominant taxa and environmental feedbacks may contribute to winner identity. We therefore clarified the implication of the Fig.~5 analysis, which distinguishes a Nutr$+$ top-down regime from a more emergent multi-species regime in Base, and integrated it with the pH-pairing analysis. Results \S2.5 now states: \mschange{``Together, the dominant-species and pH analyses suggest that Dominance in Nutr$+$ may reflect a top-down route to community-level selection, in which dominant taxa and associated environmental feedbacks help determine winner identity; however, pH contrast alone does not account for the full Dominance pattern, and Dominance in Base is less directly explained by the measured dominant-species and pH effects tested here.''} The supporting pH-pairing analysis showed that acid--alk pairings were not significantly enriched for Dominance relative to pairings in which both parental communities had similar pH in either Base or Nutr$+$ (Fisher's exact tests: Base $p = 0.49$, Nutr$+$ $p = 0.47$; Response Fig.~R1-2), although acidic parents won most Nutr$+$ acid--alk events (38/44, 86.4\%; binomial $p = 9.4 \times 10^{-7}$).

We revised the Discussion paragraph defining community-level selection accordingly. It now states, in part: \mschange{``We define community-level selection operationally as origin-correlated persistence during coalescence, rather than as a claim about one exclusive biochemical mechanism. This is a phenomenological description of the outcome, not a mechanistic attribution on its own. Under this definition, shared environmental tolerances, correlated traits, and environmental filtering by dominant pH-modifying taxa are alternative mechanistic routes that can produce the same outcome-level pattern when they make parental-community origin predictive of persistence. Our scope here is to define and test community-level selection as an outcome-level pattern, while future work can directly disentangle how these mechanisms contribute across different coalescence regimes.''}

To test environmental feedback explicitly, we ran a Ratzke--Gore-style pH-feedback model (Ratzke and Gore, 2018), which produces asymmetric replacement over an intermediate range of feedback strength (Response Fig.~R2-2). We agree that pH-mediated environmental modification is a plausible mechanism for coalescence outcomes, and we do not claim that the effective-interaction gLV is the unique explanatory framework. One observation from this comparison is that the pH-only model rarely produces Restructuring outcomes ($\sim 1\%$ across feedback strengths; 3/297 simulated events), whereas Restructuring is observed at substantially higher rates both in our experiments ($\sim 8\text{--}31\%$ across nutrient conditions) and in the effective-interaction gLV ($\sim 22\%$ across $\mu$ values). As discussed in R2-1, our gLV interaction parameter is treated as an effective coupling that can incorporate environmental-feedback contributions including pH modification; under this reading, pH-feedback and effective-interaction-strength descriptions are complementary rather than mutually exclusive.

\begin{figure}[H]
\centering
\includegraphics[width=0.88\textwidth]{internal_Q5_phase_pH.pdf}
\caption*{\textbf{Response Fig.~R2-2.} \textbf{Outcome distribution of a Ratzke--Gore-style pH-feedback model across feedback strengths.} Per-event outcomes in the similarity plane (left, three feedback strengths) and stacked outcome fractions (right) for the pH-feedback model. The model produces asymmetric replacement (Dominance) at moderate-to-strong feedback, supporting pH-mediated environmental modification as a plausible mechanism contributing to coalescence outcomes. Restructuring outcomes are rare across all feedback strengths ($\sim 1\%$; 3/297 simulated events), in contrast to the substantial Restructuring fractions observed experimentally ($\sim 8\text{--}31\%$) and in the effective-interaction gLV ($\sim 22\%$).}
\end{figure}

In addition, several related comments raise alternative theoretical models and complementary controls, including biomass effects, pH-mediated filtering, dominant-species circularity, pool size, richness, and facilitation, which are addressed in R1-1 to R1-4 and R3-2 to R3-4 and consolidated in Supplementary Note~5 (``Alternative models and controls for community-level selection''). The passive-retention alternatives are most relevant here: the denser parental community does not generally win, and we did not detect a significant overall effect of initial pool size on Dominance frequency (R1-1, R1-4).

\reviewercomment{Relatedly, we suggest clarifying the ecological meaning of ``failed invasion.'' At present it is primarily used as a proxy for interaction strength, but it is also naturally interpretable as invasion resistance (i.e.\ whether a rare invader can establish). Framing this more explicitly in terms of invasion fitness would help connect the experimental results to ecological theory and may provide additional insight into the observed patterns.}

We thank the reviewer for this important clarification. We agree that failed invasion is naturally connected to invasion resistance and invasion fitness. We therefore revised the manuscript to describe the pairwise assay as an operational measure of invasion resistance and net interaction intensity across nutrient conditions. In Fig.~4B, we retained the y-axis as the measured quantity, ``frequency of failed invasions,'' while the panel framing identifies this as an invasion-resistance analysis. Results \S2.4 now states: \mschange{``To operationally quantify its effect, we measured failed invasion frequency using pairwise invasion assays among the 12 most abundant isolates (95:5 initial frequency; Methods, Supplementary Figs.~5--7). The fraction of failed invasions increased monotonically with nutrient supply (Nutr$-$: $2 \pm 1\%$; Base: $33 \pm 4\%$; Nutr$+$: $48 \pm 4\%$; mean $\pm$ s.e.m.; Fig.~4B), indicating that invasion resistance increases across the nutrient gradient.''}

The Methods now states that \mschange{``the fraction of failed invasions served as an empirical proxy for net interaction intensity and invasion resistance.''} We distinguish this two-species assay from the pairwise selection-correlation metric in R2-4.

% ============================================
\subsection*{R2-3. Continuous measures alongside categorical outcomes}
\reviewercomment{The classification into Dominance, Mixture, and Restructuring is useful and provides a clear framework for summarising outcomes. However, it depends on thresholding a continuous similarity measure.

As noted in the broader literature [Mansour et al., 2018], coalescence outcomes often lie along a continuum rather than discrete categories.

We suggest that the authors:
\begin{itemize}
\item present continuous similarity measures alongside categorical outcomes
\end{itemize}}

\responselabel
We agree that the classification (Dominance, Mixture, Restructuring) depends on boundaries in a continuous similarity space. PDI was already introduced in Figs.~3 and 4 as a continuous measure of how strongly the coalesced community is directionally weighted toward one parental community versus the other. To make the thresholding explicit, we now present PDI together with asymmetry and retention magnitude as continuous auxiliary metrics underlying the discrete outcome classification, as shown in Response Fig.~R2-3, and have added these analyses to the Supplementary Information as Supplementary Figs.~29--31. These measures show where each coalescence event lies before category boundaries are applied, while the categorical labels provide the compact summary used in the main analyses.

To verify how sensitive our main nutrient-dependent result is to the classification-boundary choices, we varied the PDI and retention boundaries around the manuscript baseline (Response Fig.~R2-4). Across a broad range of boundary choices, the Dominance fraction remains ordered Nutr$-$ $<$ Base $<$ Nutr$+$, supporting the robustness of the nutrient-dependent trend. We also added this boundary-sensitivity analysis to the Supplementary Information as Supplementary Fig.~32 and refer to it in Supplementary Note~1.

\begin{figure}[H]
\centering
\includegraphics[width=0.9\textwidth]{marginal_distributions_by_medium.pdf}
\caption*{\textbf{Response Fig.~R2-3.} \textbf{Continuous similarity measures track the same nutrient-dependent shift as the categorical classifications.} Rows show Nutr$-$, Base, and Nutr$+$; columns show Parental Dominance Index (PDI), where values near 0 or 1 indicate strong asymmetry toward one parental community and values near 0.5 indicate balanced parental contributions; asymmetry magnitude $y$, a direction-independent measure of how unequal the two parental contributions are; and squared retention magnitude $r^2 = u^2 + v^2$, which measures how much of the coalesced composition is explained by the parental communities.}
\end{figure}

\begin{figure}[H]
\centering
\includegraphics[width=0.95\textwidth]{boundary_sensitivity_by_medium.pdf}
\caption*{\textbf{Response Fig.~R2-4.} \textbf{Boundary-sensitivity analysis supports the nutrient-dependent Dominance trend.} (A) Dominance fractions after varying the PDI boundary while holding the retention boundary at $r^2 > 0.5$. (B) Dominance fractions after varying the retention boundary while holding the PDI boundary at the manuscript baseline. Dashed lines mark the manuscript baseline. This sensitivity analysis is descriptive.}
\end{figure}

\reviewercomment{In addition, the distinction between mixture and restructuring is not entirely straightforward biologically, particularly because abundance changes contribute to classification.

We suggest that the authors:
\begin{itemize}
\item clarify the biological interpretation of ``restructuring''
\end{itemize}}

We clarify that Restructuring denotes low parental retention: the coalesced community is not well explained by either parental community alone or by a non-negative linear combination of the two parental communities. In this interpretation, Restructuring can arise when cross-community species pairings that were not present during parental assembly generate a new stable composition, especially in complex multi-species communities where abundance shifts and higher-order interactions make the final state poorly described by the parental-community basis.

Results \S2.1 (``Community-level selection is prevalent in microbial coalescence'') now includes the concise Restructuring definition: \mschange{``Restructuring denotes low parental retention: the coalesced community is not well explained by either parental community alone or by a balanced linear combination of the two parental communities, and may reflect a new stable composition produced by previously unseen cross-community species interactions.''}

\reviewercomment{We also find the comparison between dot product and Jaccard metrics in Extended Data Fig.~2 potentially very informative. The divergence between these metrics in low-nutrient environments may itself provide evidence that species-level processes dominate in that regime, and this could be discussed more explicitly.}

We agree that this comparison is informative. Extended Data Fig.~2 was generated for the Base-medium outcomes, so we discuss the metric divergence in that context. Consistent with our response to Reviewer~1, we now frame the result as a difference in what each metric measures rather than as a broad robustness claim. This pattern is consistent with the fact that the metrics emphasize different aspects of community composition: relative-abundance metrics are more sensitive to quantitative dominance in relative abundance, whereas Jaccard reflects presence/absence retention and Jensen--Shannon evaluates divergence across the full abundance distribution. Thus, a coalesced community can be skewed toward one parental community in abundance space while still retaining taxa from both parental communities, or while redistributing subdominant taxa in ways that make it less skewed toward one parental community under Jaccard or Jensen--Shannon.

Results \S2.1 now narrows the main-text robustness statement to the two alternative metrics that recovered the same pattern: \mschange{``This pattern of Dominance as the most frequent outcome was robust under similarity metric choices including Euclidean distance and Bray--Curtis dissimilarity (Extended Data Fig.~2).''} Extended Data Fig.~2 caption now carries the metric-divergence interpretation: \mschange{``Dominance is recovered by abundance-weighted parental-similarity metrics (vector decomposition, Euclidean distance, and Bray--Curtis), but not by Jaccard or Jensen--Shannon. This divergence reflects the different biological questions asked by each metric: abundance-weighted metrics quantify whether the coalesced community quantitatively resembles one parent more than the other, whereas Jaccard asks whether taxa from each parent are retained and Jensen--Shannon measures divergence of the full abundance distribution.''}

% ============================================
\subsection*{R2-4. Interpretation of pairwise selection correlation}

\reviewercomment{The pairwise ``selection correlation'' metric is an interesting attempt to quantify lineage-level coherence, but its interpretation remains somewhat unclear.

It appears closely related to co-occurrence patterns, and it is not fully clear whether the observed correlations reflect ecological interactions, shared environmental responses, or methodological effects. We find this analysis potentially insightful, but its conceptual interpretation would benefit from clarification.}

\responselabel
We agree that the pairwise selection correlation needed a clearer interpretation. Operationally, it tests, for each coalescence event, whether species pairs from the same parental community share survival outcomes more strongly than species pairs from different parental communities. The comparison is evaluated against a within-event null in which parental-origin labels are shuffled. This controls for generic shared-fate tendencies within an event, so the measured enrichment specifically asks whether shared fate is associated with parental-community origin. We use this enrichment as evidence for community-level selection, while recognizing that the metric does not identify the causal route; direct competition, shared environmental tolerance, and environmental modification can all generate such structure.

Results \S2.2 now defines the metric directly: \mschange{``We quantified this effect using pairwise selection correlation, which tests across coalescence events whether species pairs from the same parental community share survival outcomes (both persist or both go extinct) more strongly than species pairs from different parental communities (Supplementary Information; Fig.~2D).''}

\reviewercomment{We suggest adding a short paragraph linking this metric more explicitly to invasion fitness in a gLV framework. In that context, invasion fitness corresponds to the growth rate of a species when rare in a resident community. The pairwise invasion assays can then be interpreted as empirical approximations of invasion fitness in two-species systems, while coalescence outcomes reflect correlations in invasion success across many species.

This would provide a unifying theoretical framework linking pairwise assays, community coalescence, and the notion of community-level selection.}

\responselabel
We agree that this point needs clarification, especially because ``pairwise'' refers to two different analyses in the manuscript. As discussed in R2-2, the 95:5 pairwise invasion assay measures growth dynamics and ecological outcomes in two-species cocultures, and we use it primarily as an empirical approximation of invasion resistance and net interaction strength under each nutrient condition. By contrast, the pairwise selection-correlation metric is computed from community coalescence outcomes. It quantifies whether assembly history creates correlated species fates across species pairs: specifically, whether species pairs from the same parental community show more strongly coupled survival or extinction than species pairs drawn from different parental communities. In the gLV framework, this distinction corresponds to two scales of invasion fitness: the pairwise assays approximate rare-invader establishment in two-species resident backgrounds, whereas the selection-correlation metric measures whether persistence outcomes are correlated across many species invading or resisting invasion in an assembled community background. We added this conceptual bridge to Supplementary Note~2 and added an auxiliary gLV analysis showing that excess same-parent shared fate increases with the interaction-strength parameter $\mu$ (Pearson $r = 0.870$, $p = 3.2 \times 10^{-8}$; Supplementary Fig.~28).

Finally, we acknowledge the limitations and scope of this interpretation more explicitly in the Discussion: \mschange{``We define community-level selection operationally as origin-correlated persistence during coalescence, rather than as a claim about one exclusive biochemical mechanism.''} \mschange{``Accordingly, our evidence for community-level selection is operational and outcome-based: we assess it from parental-community similarity after coalescence and pairwise species-fate correlation, rather than directly measuring the mechanistic selection pressures that generate these outcomes.''} \mschange{``Thus, our scope here is to define and test community-level selection as an outcome-level pattern, while future work can directly disentangle how different mechanisms contribute across coalescence regimes.''}

% ============================================
\subsection*{R2-5. Pre-selection of natural communities}

\reviewercomment{The inclusion of natural communities is an important strength of the study. However, the use of a common stabilisation phase in defined media may introduce pre-selection effects that reduce ecological heterogeneity across communities.

In particular, incubating diverse natural communities in simplified media is known to drive convergence toward a limited set of functional guilds [Goldford et al., 2018]. As a result, the ``natural'' communities used here may already be filtered to resemble the synthetic communities, potentially contributing to the similarity of results across systems.

We therefore suggest that the authors:
\begin{itemize}
\item discuss the potential effects of this pre-selection more explicitly
\item clarify the extent to which taxonomic or functional convergence occurs during the stabilisation phase
\end{itemize}

This would help contextualise the generality of the natural-community results.}

\responselabel
We agree this is an important caveat. As the reviewer notes, defined-media stabilization can drive taxonomic or functional convergence (Goldford et al., 2018), and we cannot directly test this here because the present experiments collected 16S community profiles rather than metagenomic, metabolomic, or trait-resolved measurements. We therefore narrow our claim: the natural sample-derived experiment shows that the nutrient-dependent increase in Dominance also occurs in taxonomically richer, laboratory-stabilized communities derived from natural samples, not that unmanipulated natural ecosystems follow the same rules.

Within the available data, we asked whether stabilized communities still carried source-specific taxonomic signal. Using same-source biological replicates as a comparison baseline, different-source parental communities remained substantially less similar at the ASV level (Jaccard similarity 1.7- to 4.1-fold lower than same-source replicates across the three media; Response Fig.~R2-5b; Supplementary Fig.~44b), with a comparable pattern at the Genus level (Response Fig.~R2-5e; Supplementary Fig.~44e). This does not rule out the Goldford-style scenario in which each source converges to its own source-specific attractor during stabilization; only a pre-enrichment baseline could test that directly, and we do not have one. The data do, however, exclude the stronger version of pre-selection in which all sources collapse to a single shared taxonomic endpoint. Coalesced communities also retained parent-specific taxonomic signal, remaining closer to their own parents than to unrelated same-medium parents (all ASV/Genus tests $p<10^{-6}$; Response Fig.~R2-5c,f; Supplementary Fig.~44c,f).

We therefore limit the claim to nutrient-dependent Dominance in taxonomically richer laboratory-stabilized communities derived from natural samples, while noting that taxonomic or functional pre-selection during the seven-cycle stabilization may affect generalization to unmanipulated natural ecosystems. This revision also addresses the related Reviewer~3 concern that natural sample-derived communities are still model systems and require cautious claims (R3-4).

\begin{figure}[H]
\centering
\includegraphics[width=0.95\textwidth]{Fig_R2_5_natural_taxonomic_distinctness.pdf}
\caption*{\textbf{Response Fig.~R2-5.} \textbf{Natural sample-derived communities retain ASV- and Genus-level distinctness after stabilization and coalescence.} \textbf{a,d}, ASV- and Genus-level richness by medium. \textbf{b,e}, Same-source stabilized parental replicates are more similar than different-source parental communities, indicating retained source-specific taxonomic signal after stabilization. \textbf{c,f}, Coalesced communities remain more similar to their own parental communities than to unrelated same-medium parental communities. Analyses use ASVs above the 0.1\% relative-abundance threshold and post-stabilization 16S profiles; they test for complete post-stabilization taxonomic collapse and retained parent-specific signal, but not pre-to-post enrichment convergence or functional convergence.}
\end{figure}

We revised the manuscript to frame this experiment as applying to natural sample-derived communities established through laboratory stabilization, rather than to unfiltered natural ecosystems. We also added the same-source/different-source taxonomic-distinctness analysis to the Supplementary Information as Supplementary Fig.~44, and added a separate paragraph to Results \S2.6 to present pre-selection during stabilization as an alternative explanation for similarity between synthetic and natural sample-derived outcomes:
\mschange{``These experiments tested whether nutrient-dependent Dominance is also observed after laboratory assembly of taxonomically richer communities from natural samples. An alternative explanation is that the seven-cycle enrichment phase imposed laboratory pre-selection before coalescence. Under this hypothesis, communities that began from natural samples could have converged toward taxa or functional strategies favored by the defined media, making the subsequent coalescence outcomes resemble those of synthetic communities. The higher post-stabilization richness and retained same-source versus different-source taxonomic structure argue against complete ASV-level convergence (Supplementary Figs.~22--24 and 44), but the current data cannot determine how much of the similarity between synthetic and natural sample-derived outcomes reflects intrinsic natural-community dynamics versus pre-selection during stabilization (Supplementary Note~6).''}
We also added Supplementary Note~6 to state explicitly that pre-selection may occur and that the current data cannot quantify taxonomic convergence during stabilization or functional convergence.


% ============================================
\subsection*{R2-6. Frame gLV as phenomenological}

\reviewercomment{The gLV model successfully reproduces the qualitative transition in coalescence outcomes and serves as a useful minimal framework. However, it is limited to purely competitive interactions and does not explicitly capture environmentally mediated effects such as pH modification.

We therefore suggest framing the model more explicitly as a phenomenological rather than mechanistic description of the system.}

\responselabel
We agree that the gLV model should be framed as a phenomenological competitive baseline rather than as a mechanistic description of the biochemical interactions operating in the experiments. Results \S2.2 (``Theoretical model with random interactions reproduces community-level selection'') now states: \mschange{``The gLV model serves as a phenomenological framework to explore how the statistical properties of interaction coefficients shape coalescence outcomes, rather than as a mechanistic model of the specific biochemical interactions (e.g., pH modification, metabolic cross-feeding) underlying our experimental observations.''}

\reviewercomment{In addition, while the robustness analysis across interaction coefficient distributions is appreciated, a brief discussion of the biological interpretation of these distributions would be helpful. In particular, clarifying how the mean interaction strength parameter ($\mu$) relates to underlying ecological processes would improve interpretability.}

\responselabel
In the model, $\mu$ parameterizes the sampling range for competitive interaction coefficients in the gLV model. Larger $\mu$ makes the competitive coefficient ensemble stronger and more heterogeneous, moving the system away from near-neutral dynamics and creating stronger mismatches between species that assembled together and outsiders. This lets assembly filter species into mutually compatible groups, the ``cohesion without cooperation'' route to Dominance discussed in R3-4.

We decoupled these scale and heterogeneity effects in Supplementary Note~3 and report the corresponding simulation result in R3-4 (Response Fig.~R3-4f). Both effects contribute to Dominance, and neither alone is sufficient. This is consistent with the Hu et al.\ framework, in which $\mu$ couples them through the uniform distribution $U(0, 2\mu)$.

We also revised the Results, Methods, and Supplementary Methods to avoid presenting $\mu$ as simply a mean interaction strength and to define it instead as the sampling-range parameter for competitive interaction coefficients. Results \S2.2 now states: \mschange{``We fixed $r_i = 1$ and drew off-diagonal competition coefficients from a uniform distribution $\mathbb{U}(0, 2\mu)$, so $\mu$ controls the width of the interaction-coefficient sampling distribution.''} Methods now states: \mschange{``We fixed growth rates to 1 and self-interaction coefficients $\alpha_{ii} = 1$, and drew off-diagonal competition coefficients from a uniform distribution $\mathbb{U}(0, 2\mu)$, so $\mu$ sets the range of sampled competitive interaction coefficients within this phenomenological model.''} Supplementary Methods further defines $\alpha_{ij}$ as an effective per-capita inhibitory term and states that the baseline $\alpha_{ij}\geq 0$ model captures competitive growth inhibition rather than facilitation.


% ============================================
\subsection*{R2 minor comments}

\reviewercomment{Improve the visualisation of pairwise correlation results, as the separation between groups is not visually clear.}

\responselabel
We agree that the previous visualization could be clearer. We revised the Fig.~2D caption in Results \S2.2 to define each visual element explicitly:
\mschange{``Individual dots show per-event pairwise selection correlations; in the simulation panel, 100 stratified events per group are displayed for legibility. Squares with error bars show mean $\pm$ s.e.m.\ across all coalescence events ($n = 1{,}200$ per group in simulations and all valid experimental events in the experimental panel). Gray horizontal lines indicate the random selection baseline from a null model in which species origin labels are shuffled (see Supplementary Information).''}

\reviewercomment{Specify how pH was measured (e.g.\ continuously or at endpoints), and whether variability across replicates was assessed.}

\responselabel
We clarified the pH-measurement description in Supplementary Methods, ``Optical Density Measurements and Monoculture Characterization,'' which now specifies the timing and replicate treatment:
\mschange{``Community-level pH measurements were endpoint measurements on culture supernatant rather than continuous measurements during growth; replicate-to-replicate variability is represented by the biological replicates shown in the pH-based analyses.''}

\reviewercomment{Clarify the biological interpretation of the interaction coefficient distributions used in the theoretical model.}

\responselabel
We clarified the biological interpretation of the gLV coefficients and their distributions in Supplementary Methods, ``Lotka--Volterra Simulations.'' The revised text states:
\mschange{``Biologically, each $\alpha_{ij}$ is an effective per-capita term for the effect of species $j$ on species $i$ within the coarse-grained gLV model. It can absorb the net inhibitory consequences of direct mechanisms (e.g., resource competition, interference) and indirect mechanisms (e.g., metabolite-mediated inhibition, pH modification) at steady state, but it does not identify the underlying biochemical mechanism and the model does not dynamically represent pH or other environmental state variables.''}
We also clarify the meaning of $\mu$ and the distributional assumption:
\mschange{``Thus, in the baseline competitive ensemble, $\mu$ sets the range of sampled competitive interaction coefficients. This parameter changes the overall scale and heterogeneity of pairwise coefficients, rather than isolating a single biological mechanism.''}
\mschange{``This distribution is a phenomenological ensemble of heterogeneous effective competitive effects rather than a fitted empirical distribution of biochemical mechanisms; robustness analyses across alternative coefficient distributions and a separate mean-vs-variance sweep identify which aspects of interaction heterogeneity are required for the qualitative coalescence transition.''}

\reviewercomment{Ensure consistent terminology and notation throughout (e.g.\ ``pairwise'', ``community-level'', $\alpha_{ij}$).}

\responselabel
We standardized the terminology throughout the revised manuscript. In particular, Results \S2.2 now uses ``pairwise selection correlation'' for the lineage-level correlation metric and defines it as follows:
\mschange{``We quantified this effect using pairwise selection correlation, which tests across coalescence events whether species pairs from the same parental community share survival outcomes (both persist or both go extinct) more strongly than species pairs from different parental communities (Supplementary Information; Fig.~2D).''}
We use ``community-level selection'' for origin-correlated persistence at the parental-community level and use $\alpha_{ij}$ consistently for the gLV interaction coefficients.

\reviewercomment{Correct minor typographical issues, including ``generalis ability'' $\to$ ``generalisability''.}

\responselabel
We corrected ``generalis ability'' to ``generalisability'' and addressed other minor typographical issues throughout the revised manuscript.

\clearpage
% ============================================
% Response to Reviewer 1
% ============================================
\section*{Response to Reviewer 1}

We thank the reviewer for their constructive feedback. Below we address each point in detail.

% ============================================
\subsection*{R1-1. Absolute biomass as an alternative explanation for Dominance}

\reviewercomment{I found the Authors' explanation of coalescence generally convincing. However, I wonder if a simple, alternative explanation for many of the results could be that parent communities attain very different absolute densities. For example, if community A reaches 10x the overall density of community B, then even relatively rare species in community A might be initially more abundant than the dominant species in community B (because communities are combined in equal volumes, line 119). This might lead to the final community being dominated by species from A simply because species survive proportionally to their initial \emph{absolute} abundance. If I understand correctly, this scenario is different from the one tested in ED Fig.~5. High heterogeneity in absolute densities could potentially explain both the prevalence of Dominance outcomes and the difference in pairwise selection correlation between same- and cross-parent pairs. In principle, heterogeneity in absolute densities could also increase with nutrient concentration. To be clear, I have no reason to expect this pattern, but I wonder if the Authors could do more to rule it out. Since ODs were measured throughout (line 379), could those data be used to check this possibility?}

\responselabel
We thank the reviewer for raising this important alternative explanation. We tested the suggested hypothesis by asking whether parental-community OD$_{600}$ correlates with the parental dominance direction or with pairwise selection correlation.
If Dominance were primarily driven by unequal absolute biomass at mixing, the higher-OD parental community should win more frequently, and signed $\Delta\text{OD}_{A-B}$, the difference between the OD of parental community A and the OD of parental community B, should be positively associated with PDI. We observed the opposite pattern. In Base, the higher-OD parental community wins only $37\%$ of Dominance events, and the continuous relationship between signed $\Delta\text{OD}_{A-B}$ and PDI is weak. This opposite-to-prediction pattern becomes strongest in Nutr$+$, where Dominance is most frequent: the higher-OD parental community wins only $9/68 = 13.2\%$ of Dominance events ($p = 3.9 \times 10^{-10}$), and signed $\Delta\text{OD}_{A-B}$ is strongly negatively correlated with PDI ($\rho = -0.60$, $p = 4.0 \times 10^{-10}$; Response Fig.~R1-1A,B). Nutr$-$ is less informative for this directional test because both OD differences and Dominance frequency are smaller. Thus, this directional analysis argues against parental-community biomass differences as an explanation for Dominance direction.

This density-based explanation also predicts that higher-OD parental communities should show stronger within-community pairwise selection correlations. To test this, we binned parental communities into low, mid, and high OD tertiles within each medium and computed the pairwise selection correlation among species from the same parental community in each OD bin (Response Fig.~R1-1C). The cross-community correlation is not re-binned by OD and is shown as a medium-level reference. Within-community pairwise selection correlation does not show a consistent increase with parental-community OD; instead, the pattern is not consistent across media. In Nutr$+$, the slight negative OD trend mirrors the negative signed $\Delta\text{OD}_{A-B}$--PDI association in the directional analysis above, rather than the positive trend expected if higher biomass drove winning. Overall, this analysis implies that differences in parental-community biomass do not explain Dominance direction or the nutrient-dependent pattern of correlated species fates.

We added a concise note to the Results and present the primary analysis in Supplementary Note~5 and Supplementary Figs.~35--36. The Results now state: \mschange{``We next tested whether nutrient-dependent changes in parental-community properties, including biomass heterogeneity and initial pool-size effects, could explain this trend; these factors did not account for the nutrient-dependent increase in Dominance or correlated species fates (Supplementary Note~5 and Supplementary Figs.~34--36).''}

\begin{figure}[H]
\centering
\includegraphics[width=\textwidth]{Fig_R1_1A_winner_loser_OD.pdf}
\caption*{\textbf{Response Fig.~R1-1A.} \textbf{Higher parental-community OD does not preferentially predict the winner in Dominance events.} Winner OD vs loser OD scatter is shown per medium; dashed line is $y=x$. Points above the line are cases where the winning parental community had higher OD than the losing parental community. The higher-OD parental community wins only $37\%$, $37\%$, and $13\%$ of Dominance events in Nutr$-$, Base, and Nutr$+$ respectively; the strongest opposite-to-prediction pattern occurs in Nutr$+$. The displayed $p$ value is from a two-sided binomial test against the null expectation that the higher-OD parental community wins $50\%$ of Dominance events.}
\end{figure}

\begin{figure}[H]
\centering
\includegraphics[width=\textwidth]{Fig_R1_1B_OD_vs_PDI.pdf}
\caption*{\textbf{Response Fig.~R1-1B.} \textbf{Continuous OD differences do not support the prediction that the higher-OD parental community wins.} Signed $\Delta\text{OD}_{A-B}$ ($\text{OD}_{A} - \text{OD}_{B}$) vs PDI is shown per medium. PDI is the Parental Dominance Index defined in the manuscript, $\text{PDI} = (2/\pi)\arctan(u/v)$. Spearman $\rho$ is the rank correlation between signed $\Delta\text{OD}_{A-B}$ and PDI; the displayed $p$ value is from a two-sided Spearman rank-correlation test of the null hypothesis of no monotonic association. The relationship is weak in Base and strongly negative in Nutr$+$ ($\rho = -0.60$, $p = 4 \times 10^{-10}$), opposite the sign predicted if higher OD caused a parental community to win. Gray points are reflections included for symmetry.}
\end{figure}

\begin{figure}[H]
\centering
\includegraphics[width=\textwidth]{Fig_R1_1C_pairwise_corr_vs_OD.pdf}
\caption*{\textbf{Response Fig.~R1-1C.} \textbf{Higher parental-community OD does not produce a consistent increase in within-community pairwise selection correlation.} Columns show Nutr$-$, Base, and Nutr$+$, respectively. \textbf{Top row:} parental communities are grouped into low, mid, and high within-medium OD tertiles; smaller text below the group labels shows the OD$_{600}$ range of each tertile. Bars show within-community pairwise selection correlation with bootstrap 95\% CIs; the dashed gray line and shaded band show the medium-level mean cross-community correlation and its bootstrap 95\% CI. \textbf{Bottom row:} per-event parental-community observations plotted by OD and within-community pairwise selection correlation; vertical dotted lines mark the low/mid/high OD boundaries used in the top row. Spearman $\rho$ and $p$ values are shown in each panel; p-values are from two-sided Spearman rank-correlation tests of the null hypothesis of no monotonic association between OD and within-community pairwise selection correlation.}
\end{figure}

% ============================================
\subsection*{R1-2. Does pH mismatch predict Dominance?}

\reviewercomment{It seems like an even more straightforward prediction would be that Dominance should become more likely when combining parent communities with qualitatively different pH (e.g.\ alkaline and acidic) compared to cases where both parents are either acidic or alkaline. Is this the case?}

\responselabel
We thank the reviewer for suggesting this direct test. Response Fig.~R1-2 shows that Dominance frequency is not significantly different between acid--alk and same-pH parental-community pairs. For this frequency-level analysis, we compared acid--alk pairs with pooled same-pH pairs (acid--acid or alk--alk) within each medium. Nutr$-$ is excluded from this analysis because parental communities show little pH change in this condition. Compared with pooled same-pH pairs, acid--alk pairs showed a slightly higher Dominance frequency in Nutr$+$ ($79.5\%$, 35/44, versus $71.7\%$, 33/46; OR $= 1.53$, Fisher $p = 0.47$), but a slightly lower Dominance frequency in Base ($61.4\%$, 27/44, versus $69.4\%$, 25/36; OR $= 0.70$, Fisher $p = 0.49$). This result implies that pH mismatch alone is not sufficient to explain Dominance frequency: even in Nutr$+$, where Dominance is common across pH-pair types, the acid--alk enrichment is not statistically distinguishable from same-pH pairs.

\reviewercomment{To support the hypothesis that the top-down regime emerges because dominant species strongly modify the pH of the ecosystem, the Authors show that the predictive power of the parent community pH increases in the Nutr+ condition when looking at coalescence between communities where one parent is acidic and one is alkaline (lines 261, ED Fig.~8).}

Yes, considering only the pH-mismatched acid--alk group amplifies the dominance-direction signal. In Nutr$+$, the lower-pH parental community dominates 72\% of all pH-pair types (binomial $p = 3.0 \times 10^{-5}$), but the acidic parental community dominates 86.4\% of acid--alk events (38/44; binomial $p = 9.4 \times 10^{-7}$; Response Fig.~R1-2). In Base, the all-pair signal is weak (54\%, binomial $p = 0.58$), and the acid--alk-only signal remains statistically insignificant (63.6\%, 28/44; binomial $p = 0.096$). Fig.~5 and the associated predictability analysis remain based on all eligible events, including both same-pH and pH-mismatched pairs. The pH-mismatch-specific analysis is reported separately in Supplementary Note~5 and Supplementary Fig.~37.

We therefore tuned the overall claim in Results \S2.5 to integrate the pH analysis with the dominant-species predictability analysis, while making clear that pH contrast alone does not explain the full Dominance pattern. The revised Results paragraph states: \mschange{``We next asked whether environmental modification through pH could contribute to the top-down regime. In our system, dominant species are often strong pH modifiers that either acidify or alkalinize the medium (Supplementary Fig.~8). In both Base and Nutr$+$ media, when these pH-modifying species become dominant within a community, they are predictive of the community's overall pH; communities dominated by acidifiers become acidic, while those dominated by alkalizers become alkaline (Supplementary Fig.~9). In coalescence events between acidic (pH $<$ 6.5) and alkaline (pH $>$ 7.5) parental communities, the acidic community dominated in only 56\% of cases in Base medium ($n = 41$), but dominated in 91\% of cases in Nutr$+$ medium ($n = 32$, Fisher's exact test $p < 0.0001$; Extended Data Fig.~8), consistent with the hypothesis that high nutrients amplify metabolic activity, thereby intensifying pH modification and excluding pH-sensitive species. A complementary pH-pairing analysis showed that pH effects do not explain the full Dominance pattern: acid--alk pairings were not significantly enriched for Dominance relative to pairings in which both parental communities had similar pH in either Base or Nutr$+$ medium (Fisher's exact tests, $p = 0.49$ and $p = 0.47$; Supplementary Fig.~37). Together, the dominant-species and pH analyses suggest that Dominance in Nutr$+$ may reflect a top-down route to community-level selection, in which dominant taxa and associated environmental feedbacks help determine winner identity; however, pH contrast alone does not account for the full Dominance pattern, and Dominance in Base is less directly explained by the measured dominant-species and pH effects tested here.''}

\begin{figure}[H]
\centering
\includegraphics[width=\textwidth]{Fig_R1_2_acidalk_per_medium.pdf}
\caption*{\textbf{Response Fig.~R1-2.} \textbf{pH contrast alone does not significantly predict Dominance frequency, but lower-pH parental communities preferentially dominate in Nutr$+$.} \textbf{(Left two panels)} Stacked class-fraction bars (Dominance / Mixture / Restructuring) for the three pH-pair types in Base and Nutr$+$; pair types are ordered Acid--Alk, Acid--Acid, Alk--Alk, with Dominance \% labeled inside each bar. Reported class-frequency p-values are from two-sided Fisher exact tests on 2$\times$2 tables comparing Acid--Alk vs pooled same-pH (Acid--Acid or Alk--Alk) events and Dominance vs non-Dominance outcomes. (All pairs) Signed asymmetry for all pH-pair types: positive = lower-pH parental community dominated, negative = higher-pH parental community dominated. (Acid--Alk pairs only) Signed asymmetry restricted to acid--alk pairs: positive = acidic parental community dominated, negative = alkaline parental community dominated. Percentages above the signed-asymmetry panels report lower-pH or acidic parental-community dominance fractions; p-values above these panels are from two-sided binomial tests against the null expectation that the lower-pH or acidic parental community dominates 50\% of events. Dots = individual events; squares = mean $\pm$ s.e.m.}
\end{figure}

% ============================================
\subsection*{R1-3. Circularity in Fig.~5C, PDI excluding dominant species}

\reviewercomment{In Fig.~5C, it seems like there is some risk of circularity if the dominant species are included in the calculation of the PDI. Does the positive correlation under Nutr+ conditions hold up if the dominant species are excluded from this calculation?}

\responselabel
We agree that this is an important circularity check. We used two removal strategies to avoid biasing the result toward the mixed-community winner. In the first, we removed the most abundant species in the coalesced community from all three community composition vectors; in the second, stricter test, we removed the most abundant species from each parental community separately before recalculating PDI. We followed the same PDI calculation and filtering criteria as in the original Fig.~5C analysis: events are included if they are classified as Mixture or Dominance before dominant species removal.

This analysis reveals that our core conclusion remains supported: dominant-species pairwise interactions predict the community-level winning direction in Nutr$+$, but not in Base. In Base, $R^2$ drops from $0.11$ to $0.00 / 0.07$ and the slope from $0.49$ to $-0.02 / 0.33$ (Original / mix-winner removed / parental-community dominants removed). In Nutr$+$, $R^2$ drops from $0.49$ to $0.24 / 0.18$ and the slope from $1.03$ to $0.64 / 0.56$ (Response Fig.~R1-3). Using independent, non-reflected events, the Spearman correlation between dominant-species pairwise score and community-level score remains significant after both removal strategies: $\rho = 0.41$ ($p = 8.2 \times 10^{-4}$) after removing the mix-winner species and $\rho = 0.42$ ($p = 1.3 \times 10^{-3}$) after removing the dominant species from both parental communities. The corresponding linear slopes are also positive and significant ($p = 1.1 \times 10^{-4}$ and $1.5 \times 10^{-4}$). Thus, the dominant species is the primary single contributor in Nutr$+$, but the Nutr$+$ trend does not disappear when dominant species are excluded. In contrast, Base does not retain a significant correlation after removal ($p = 0.94$ and $0.12$ for the two linear slopes; Spearman $p = 1.00$ and $0.21$).

We added a concise note to Results \S2.5 and present the primary analysis in Supplementary Fig.~33. The Results now state: \mschange{``The Nutr$+$ predictability trend remained significant after recalculating PDI with dominant species excluded (Spearman rank-correlation tests: $p = 8.2 \times 10^{-4}$ after removing the coalesced-community dominant species and $p = 1.3 \times 10^{-3}$ after removing the two parental-community dominant species; Supplementary Fig.~33), indicating that it is not solely a PDI circularity artifact.''}

\begin{figure}[H]
\centering
\includegraphics[width=\textwidth]{Fig_R1_3_per_medium_scatter.pdf}
\caption*{\textbf{Response Fig.~R1-3.} \textbf{Nutr$+$ predictability persists after dominant-species removal, arguing against a PDI circularity artifact.} Per-medium scatter plots are shown with no pooling between Base and Nutr$+$. Rows: Base (MN, brown) and Nutr$+$ (HN, orange). Columns: original Fig.~5C calculation; dominant species from the coalesced community removed from all three composition vectors; and dominant species from each parental community removed before recalculating PDI. Colored points are independent coalescence events; gray points are reflected counterparts plotted for symmetry. The third column is the stricter test. $R^2$ and slope annotations are computed on the reflected plotting data, as in Fig.~5C; any displayed slope p-values are from two-sided ordinary least-squares tests of slope $=0$. Spearman $\rho$ and $p$ are computed on the independent, non-reflected events; Spearman p-values are from two-sided rank-correlation tests of the null hypothesis of no monotonic association.}
\end{figure}

% ============================================
\subsection*{R1-4. Pool-size effects}

\reviewercomment{I wondered if the Authors could provide slightly more detail on the effect of the initial pool size. Is it similarly true that there is no significant effect of initial richness in the experimental communities? By eye, it looks like Dominance might be more likely for the 24 species pool, but it is hard to tell by squinting at Fig.~1E. It might also be interesting to see the survival ratio as a function of initial richness for the experimental communities.}

\responselabel
We do not detect a significant effect of pool size on Dominance frequency in the experiment, and we thank the reviewer for prompting this more direct check. Response Fig.~R1-4C shows Dominance frequency by initial pool size within each nutrient condition; the pool-size effect is not significant either when pooled across nutrient conditions ($\chi^2 = 2.24$, $p = 0.69$) or within individual nutrient conditions (Dominance vs non-Dominance tests: Nutr$-$ $p = 0.39$, Base $p = 0.82$, Nutr$+$ $p = 0.27$). Thus, although initial pool size affects realized richness and assembly survival ratio (panels A,B), these assembly effects do not translate into a significant change in Dominance frequency.

\reviewercomment{I also wondered if the Authors could look more closely at why there is no effect in the model. Perhaps by plotting the realized richness and interaction strength in parent communities as a function of initial richness?}

To address the model part of the comment, we performed a gLV pool-size ablation study at low, intermediate, and high values of the interaction-strength parameter $\mu$, varying the initial pool size from 4 to 48 species per parental community (100 replicate assemblies per $\mu$ and pool size). The model gives a consistent interpretation: at fixed $\mu$, Dominance frequency changes only modestly with pool size compared with the strong effect of $\mu$ itself (panel E), whereas the same-parental-community versus cross-parental-community pairwise-selection gap increases with pool size (panel F). Thus, pool size can modulate the strength of correlated species fates, but it does not explain the primary interaction-strength-dependent transition in Dominance frequency.

We added this analysis to Supplementary Note 5, ``Alternative models and controls for community-level selection,'' and to Supplementary Fig.~34. The new supplementary text states that pool-size variation affects assembly richness but does not explain the Dominance-frequency pattern, and that the model pool-size ablation supports the same conclusion.

\begin{figure}[H]
\centering
\includegraphics[width=\textwidth]{pool_size_analysis.pdf}
\caption*{\textbf{Response Fig.~R1-4.} \textbf{Dominance frequency is primarily associated with nutrient condition or model interaction strength, rather than initial pool size.} \textbf{A:} Realized parental richness. \textbf{B:} ASV-based assembly survival ratio, computed as realized parental ASV richness divided by inoculated species-pool size; this ratio can exceed one because the numerator is thresholded ASV richness rather than a count restricted to the inoculated isolates. \textbf{C:} Experimental Dominance fraction by medium and initial pool size. \textbf{D:} Experimental pairwise species selection grouped by initial pool size. \textbf{E:} Model Dominance fraction grouped by $\mu$. \textbf{F:} Model pairwise species selection across pool sizes $4$--$48$.}
\end{figure}

% ============================================
\subsection*{R1-5. Similarity-metric robustness claim}

\reviewercomment{The Authors write (line 134) that ``This pattern of Dominance as the most frequent outcome was robust across variants of similarity metrics''. I liked the Author's approach to measuring outcomes, and I don't advocate for changing it, but this claim seems a little misleading. Two of four alternative similarity metrics gave different results, with Dominance being the \emph{least} likely outcome under Jensen--Shannon and Jaccard.}

\responselabel
We agree that our original robustness claim was too broad and could be misleading: the Dominance pattern is recovered by Euclidean distance and Bray--Curtis dissimilarity, but not under Jaccard index or Jensen--Shannon divergence (Extended Data Fig.~2). This pattern is consistent with the fact that the metrics emphasize different aspects of community composition: abundance-weighted parental-similarity metrics such as vector decomposition, Euclidean distance, and Bray--Curtis dissimilarity are more sensitive to quantitative dominance in relative abundance, whereas Jaccard reflects presence/absence retention and Jensen--Shannon evaluates divergence across the full abundance distribution. Thus, a coalesced community can be skewed toward one parental community in abundance space while still retaining taxa from both parental communities, or while redistributing subdominant taxa in ways that make it less skewed toward one parental community under Jaccard or Jensen--Shannon.

We revised Results \S2.1 to narrow the main-text claim to the two alternative metrics that recovered the same pattern: \mschange{``This pattern of Dominance as the most frequent outcome was robust under similarity metric choices including Euclidean distance and Bray--Curtis dissimilarity (Extended Data Fig.~2).''} We also revised the Extended Data Fig.~2 caption to explicitly discuss the divergent metrics: \mschange{``The primary Dominance pattern is recovered under similarity metric choices including Euclidean distance and Bray--Curtis dissimilarity, but not under Jensen--Shannon divergence or Jaccard index, for which the coalesced community is less often measured as skewed toward one parental community. This divergence reflects the fact that the metrics emphasize different aspects of community composition: vector decomposition, Euclidean distance, and Bray--Curtis dissimilarity capture quantitative similarity to each parental community, whereas Jaccard reflects presence/absence retention and Jensen--Shannon evaluates divergence across the full abundance distribution. Thus, Jaccard and Jensen--Shannon are best interpreted as complementary metrics rather than as failed robustness checks for the abundance-weighted Dominance classification.''}


% ============================================
\subsection*{R1-6. Gray reflected points in figures}

\reviewercomment{It took me quite a while to realize that the gray points in Figs.~1E, 4C, 5C, and 6B are simply a reflection of the colored points. Please indicate this clearly, or perhaps just remove these.}

\responselabel
We retained the reflected points and updated the captions of Figs.~1E, 4C, 5C, and 6B to state explicitly that the gray points are reflections: \mschange{``Gray points represent the symmetric counterpart of each coalescence event (community B versus A).''}

We revised the captions of Figs.~1E, 4C, 5C, and 6B as quoted above.


% ============================================
\subsection*{R1-7. Interaction matrix after assembly}

\reviewercomment{In Fig.~2A, it might be helpful to show the matrix of interaction coefficients after assembly of the parent communities (with visible block structure).}

\responselabel
We decided not to add this matrix to Fig.~2D, because Fig.~2D already focuses on pairwise selection correlation and adding interaction-matrix panels would make the main figure overly dense. Instead, we added Supplementary Fig.~27 to show the requested post-assembly interaction matrices directly. The new supplementary figure shows three representative before-versus-after assembly examples, with survivor counts annotated for each parental community, and a pooled comparison showing that assembly filters surviving species into within-community groups with lower mutual competition than the unfiltered between-community coefficients at $\mu = 0.50$.

% ============================================
\subsection*{R1-8. Fig.~2D visualization}

\reviewercomment{Fig.~2D, I was somewhat confused by the relationship between points and squares (means?). Visually, the squares do not look like they are near the means of the points. Please clarify this for readers. Additionally, the gray horizontal bars are not explained in the main text unless I missed it.}

\responselabel
The previous Fig.~2D was unclear because the individual dots and summary markers represented different display levels. In the simulation panel, the dots show a stratified subset of per-event pairwise selection correlations for visual legibility, whereas the squares and error bars show the mean $\pm$ s.e.m.\ across all simulated coalescence events. This explains why the squares need not visually coincide with the apparent center of the displayed dots. We revised the Fig.~2D caption to state this distinction explicitly and to define the gray horizontal lines as the random-selection baseline from an origin-shuffle null model: \mschange{``Individual dots show per-event pairwise selection correlations; in the simulation panel, 100 stratified events per group are displayed for legibility. Squares with error bars show mean $\pm$ s.e.m.\ across all coalescence events ($n = 1{,}200$ per group in simulations and all valid experimental events in the experimental panel). Gray horizontal lines indicate the random selection baseline from a null model in which species origin labels are shuffled (see Supplementary Information).''}


% ============================================
\subsection*{R1-9. Emphasize ``cohesion without cooperation''}

\reviewercomment{Discussion, lines 342--345: I think it is worth emphasizing how interesting and perhaps surprising it is to find ``community-level cohesion without cooperation'' as the paper by M.\ Tikhonov has it. This finding will be counterintuitive to many readers, even though the mechanism is made clear.}

\responselabel
We agree with the reviewer that this is a central and counterintuitive implication of the study. We added a dedicated short paragraph to the Discussion that frames the result in these terms, placed immediately after the paragraph reconciling community-level and species-level selection across prior studies. The new paragraph emphasizes that community-level cohesion can arise without cooperative interactions, and situates this point in relation to Tikhonov's prediction and classical random Lotka--Volterra theory.

\mschange{``Notably, and perhaps counterintuitively, our results demonstrate that community-level cohesion can emerge without cooperative interactions between species, a prediction made by Tikhonov (2016) using resource-competition models and consonant with classical random Lotka--Volterra theory (May, 1972; Grilli et al., 2017). This finding is surprising because community-level behavior is often associated with cooperative mechanisms such as cross-feeding or division of labor, whereas our results indicate that assembly history and competitive interactions can also generate correlated species fates during coalescence.''}


% ============================================
\subsection*{R1-10. Means missing in Extended Data Fig.~5C}

\reviewercomment{In ED Fig.~5C, are means missing?}

\responselabel
We thank the reviewer for pointing this out. The means are shown in Extended Data Fig.~5C as square markers with error bars, but they were not sufficiently clear in the previous version. We revised Extended Data Fig.~5 so that the mean markers and error bars are more visible, and the caption now states that individual dots show per-event correlations whereas squares with error bars show mean $\pm$ s.e.m.


% ============================================
\subsection*{R1-11. Incorrect Extended Data reference in the SI}

\reviewercomment{On p.~9 in the SI, I believe the sentence ``Simulations were performed with parental community sizes ranging from 4 to 48 species per community (Extended Data Fig.~5)'' should refer to ED Fig.~4.}

\responselabel
We thank the reviewer for catching this error and corrected the SI reference.


\end{document}
